\section{Discussion}
\label{sec:discussion}

\paragraph{What the results show.} At a real Mars site, the constraints, not the optimizer, decide most plans. A naive slope-masked planner broke documented constraints almost every time, mostly ones slope never sees, and three quarters of its routes contained steps that cell-slope masks miss. A greedy planner with the same checks matched exact selection whenever time was plentiful. The value of a crew-side planner therefore lies first in checking constraints and explaining them, and second in optimizing science, where exact selection helped when time was short. The staging result shows why checks must run along the trajectory: per-cell rules make the habitat and the rover look identical, while heat rests and accumulated radio gaps separate them. Suit thermal capacity and communications coverage also interact. From the habitat, raising the heat budget only exposed the radio limit, while from the rover a 530\,W budget added a seventh stop.

\paragraph{Why the architecture fits Mars.} Earth cannot answer in time, the helmet cannot afford dense fields, and the link will drop behind terrain. A signed package with a conservative coarse grid lets the navigator make the decisions that cannot wait, such as whether the crewmember is still within walk-back and shelter limits and which way is home, with no link and no habitat computation. The same split would serve lunar operations where line of sight to the base is intermittent.

\paragraph{Limitations.} No suited person has walked a \Marswalk plan, no EVA operator has used the console, and the explanations are untested with users. The walker in our tests is a simulator that uses the engine's own power model. The simulator shares models with the planner, so it catches search errors, not model errors. The metabolic model extrapolates to Mars gravity and under-predicts level walking. Its standing cost (474\,W) exceeds both the heat budget and the 250\,W rest power we assume, and the 30-minute window, rest power and initial heat state were not varied. Walk-back and shelter fields ignore heat rests, which presumes that contingency returns may exceed the nominal heat budget. Oxygen is a time budget, not a consumables model. Line of sight ignores curvature, Fresnel clearance and link budget, and the rover relay is assumed. Three constraint values are team parameters, and the station values are a team heuristic. Station selection is exact only over precomputed legs and up to 15 candidates, and exactness was tested against brute force only up to seven stations. The link is a local WebSocket, not a radio. Only route integrity is protected; the link is unencrypted and consoles are unauthenticated. The AR navigator's positioning is simulated, and the suit hardware is a concept. The impact and sensitivity studies use one site, one habitat position and single end-to-end runs.

\paragraph{Next steps.} The most valuable next steps are human-in-the-loop: replaying analog EVA data through the engine and putting the console in front of EVA operators. On the modelling side, the obvious checks are a gravity-independent floor for walking energy and comparisons with NASA's own regression for the suited data and with Apollo EVA metabolic rates \citep{waligora1975}. Further steps are heat-aware contingency fields, radio-aware routing and a sweep over habitat positions. On the hardware side, the concept study proposes replaying recorded boot-IMU and camera data through the estimator, an analog walk with surveyed checkpoints, worst-case reconnection tests and helmet readability trials. The pipeline accepts any digital elevation model, so lunar sites are a natural second domain.

\section{Conclusion}
\label{sec:conclusion}
Mars crews will plan and carry out many of their own EVAs, with Earth up to 22 light-minutes away and out of reach near conjunction. \Marswalk shows that a crew-side planner for that job can be built from public NASA data, published physiology and well-understood algorithms. It plans against twelve sourced constraints, re-checks every plan in a separate simulator, explains each verdict, and hands the astronaut a signed route that keeps guiding without a link or GPS. Its models are compared with measured NASA data, and their errors are reported in \cref{sec:eval-models}. At Gale it yields a concrete staging result: relaxing either the heat limit or the radio-gap limit alone leaves the habitat walk infeasible, while a nearby pressurized rover makes a six-stop walk feasible. Suited people, real radios and EVA operators come next; the architecture, code and evidence are open for others to test and improve.
